\section{Results}
\subsection{Complete cohort and account coverage}
The website and API identifier sets matched exactly: 917 model repositories,
with no duplicate identifiers or missing pagination results. The Hub
classified 409 of their 495 publisher accounts as users and 86 as
organisations, accounting for 748 and 169 model repositories respectively.
These platform account types are not findings about funding independence.
Recorded parameter counts ranged from 2 to 998,008,576; the very small
counts are retained as catalogue observations rather than interpreted as
complete language-model sizes.

Table~\ref{tab:coverage} gives the complete account-level outcomes.
The 245 unresolved accounts constitute 49.5\% of publishers and account
for 419 of the 917 repositories. All remain in the released table.
The mixed category includes both international collaborations and
conflicting personal-location evidence. The non-UK category also includes
clearly labelled organisational locations or professional affiliations;
it is not a count of 235 individually verified non-UK residents.

\begin{table}[htbp]
\centering
\caption{Location evidence for all publishing accounts. Model counts in
this table inherit the publisher's category; model-specific UK connections
are reported separately below.}
\label{tab:coverage}
\begin{tabular}{lrr}
\toprule
Account evidence category & Accounts & Model repositories \\
\midrule
Explicit UK public location & 8 & 8 \\
UK affiliation only & 1 & 1 \\
Mixed or conflicting locations & 6 & 30 \\
Non-UK geographic evidence & 235 & 459 \\
Unresolved & 245 & 419 \\
\midrule
Total & 495 & 917 \\
\bottomrule
\end{tabular}
\end{table}

\subsection{Eight individually published UK findings}
The eight accounts with explicit UK personal locations each contributed
one repository to the cohort. Table~\ref{tab:individuals} lists every case.
The location links point to the public identity evidence; the references
support the model descriptions. Current location is not automatically the
location of an earlier training run.

\begingroup
\small
\setlength{\tabcolsep}{5pt}
\begin{longtable}{@{}>{\raggedright\arraybackslash}p{0.26\linewidth}>{\raggedright\arraybackslash}p{0.18\linewidth}>{\raggedright\arraybackslash}p{0.49\linewidth}@{}}
\caption{All eight individual publisher accounts with explicit UK locations.}
\label{tab:individuals}\\
\toprule
Researcher and release & Public location & Documented work \\
\midrule
\endfirsthead
\toprule Researcher and release & Public location & Documented work \\
\midrule
\endhead
\bottomrule
\endfoot
Phillip C. O'Brien (Eden Eldith)\newline
\path{edeneldith/WiggleGPT}
& \href{https://garden-backend-three.vercel.app/}{Gosport}
& 123.7M parameters; pretrained from scratch on OpenWebText, with a
learnable oscillating MLP activation. One consumer GPU at a time, first
RTX 3070 then RTX 5060 Ti. Explicitly describes independent work without
institutional funding \cite{wigglecard}. \\[0.6em]
Ahmad Bader\newline\path{Ahmadbader17/tiny-gpt}
& \href{https://github.com/Ahmadbader17}{Cambridge}
& 5.85M GPT-2-style model implemented and trained from scratch on a laptop
CPU. TinyStories, 77.8M training tokens, 7.9 hours on an Intel Core Ultra 9
185H. Learning/demo use; no project-funding statement \cite{tinygpt}. \\[0.6em]
Sam Schorb\newline\path{iimaginary/schorbGPT-medium}
& \href{https://github.com/sam-schorb}{Sheffield}
& GPT-2 architecture and tokenizer, described as a web-text base model
without fine-tuning. Initialisation history, hardware and funding are
not specified in the reviewed sources \cite{schorb}. \\[0.6em]
Sam Dobson\newline\path{sdobson/tinystories-llama-15m}
& \href{https://github.com/samdobson}{UK}
& 15M Llama-style TinyStories model, with an implementation adapted from
\texttt{llama2.c} and 100,000 reported iterations. Initialisation and
actual hardware are not explicitly stated \cite{dobson}. \\[0.6em]
Sripaad Srinivasan\newline\path{Isotonic/bullet-points-generator}
& \href{https://github.com/sripaad}{London}
& T5 conditional bullet-point generator. The autogenerated card claims
from-scratch training on 35,000 source/target rows; its initialisation
history is not otherwise established \cite{isotonic}. \\[0.6em]
Rajesh Kanna V\newline\path{XBOT-RK/distilgpt2-wiki-qa}
& \href{https://github.com/vrajeshtrichy}{UK}
& Explicitly fine-tuned from DistilGPT2 for Wiki-QA. The Base-only
catalogue selection did not exclude this derivative \cite{xbot}. \\[0.6em]
Sally Sonia Simmons\newline\path{SallySims/equibert-generator}
& \href{https://github.com/SallySims}{UK}
& T5-base fine-tuned for conditional text generation. Separate encoder
architecture descriptions elsewhere in the card are not attributed to
this generator \cite{sally}. \\[0.6em]
Abdullatif K\"oksal\newline\path{akoksal/LongForm-OPT-125M}
& \href{https://akoksal.com/}{London}
& OPT instruction fine-tuning using LongForm data, produced through
academic collaboration. The publisher currently works at Google DeepMind
in London; this does not establish the earlier training site \cite{longform}. \\
\end{longtable}
\endgroup

Ahmad Bader's release and WiggleGPT provide direct, documented examples
of training from scratch on consumer hardware. They contribute different
kinds of evidence. The former is a small standard-architecture educational
model trained on a laptop CPU. WiggleGPT is an architecture experiment
with a learned oscillating activation, a 600,000-step pretraining record,
released checkpoints and a stated absence of institutional funding.
Its recorded pretraining ran from 23 October to 12 November 2025, with
later release updates and task adaptations. This dated research record
is substantive capability evidence even though its creator has no
registered AI company.

The remaining personal releases demonstrate why implementation,
pretraining, fine-tuning, architecture and funding should be separate
fields. They contribute to the discovery record without being forced
into a single category of ``independent from-scratch LLM''.

\subsection{UK connections within collective releases}
Following named researchers beyond publisher addresses added eight
repositories with a person's public UK location. These comprise three
Outpost Bio Waypoint models through coauthor Neythen Treloar's London
location; three PolyPythias 160M seed releases through joint first author
Pietro Lesci's public Cambridge profiles; BLOOMZ-560M through coauthor
Samuel Albanie in London; and BLOOM-560M-intermediate through dataset
and applications contributor Daniel van Strien's UK profile
\cite{waypoint,polypythias,bloomz,bloom}.

The resulting 16 repositories span 11 publisher accounts and identify
12 named people, counting the author's public names as one person.
Fifteen repositories have an identified publisher/developer or model-paper
coauthor with UK location evidence. The sixteenth is retained explicitly
as a dataset/applications contribution, not evidence that van Strien
personally trained that checkpoint. Coauthorship does not establish who
operated the training run.

The Waypoint releases model microbiome taxonomic sequences, not ordinary
English prose. Their English metadata tag admitted them to the original
query; the task-specific evidence identifies the difference. The remaining
13 of these 16 repositories concern natural-language generation, including
task-specific and multilingual releases. Collective or company research
is retained without describing it as unaffiliated solo work.

Thirteen further repositories have UK institutional-affiliation evidence:
Augusto Corr\^ea's Oxford-linked SmolLM fine-tune; GPT-Neo-125M through
named author Sid Black's UK institutional role; and eleven original
Pythia-suite releases through Herbie Bradley's Cambridge affiliation
\cite{correa,gptneo,pythia}. Bradley's documented contribution was WinoBias
analysis, and his current personal site locates him in San Francisco.
These rows are kept in an affiliation view rather than assigned current
UK residence. Seven other reviewed collective releases had no established
model-specific UK link, despite UK participation elsewhere in the larger
organisation or model family.

\subsection{A location conflict and partial training record}
\path{AiRukua/BabyKilat} describes a custom hybrid model trained from
scratch, released at 6,000 of 50,000 planned steps \cite{babykilat}.
Its author's linked GitHub profile says Southampton, whereas LinkedIn
identifies Ambon, Indonesia and records Southampton study. The conflicting
location evidence is retained. The card gives estimated training times
for several GPUs, which do not identify the device actually used.
This case is visible in the broader 30-repository UK-connected view,
with both the location conflict and partial-training status explicit.

\subsection{Catalogue filters and substantive evidence}
Thirty-seven model-level assessments cover the personal findings,
collective cases and unresolved related variants. Twenty-three of all
917 catalogue matches contain a nonempty declared \texttt{base\_model}
field. Directly inspected cards include fine-tuned checkpoints despite
the Base-only filter. Some model cards describe a different task from
the catalogue task; others are untrained or debugging artefacts.

Parameter metadata also requires interpretation. For example,
\path{ml-resources/MobileMoE-S-QAT} appears in the under-one-billion
cohort, while its card describes 272M active parameters and 1.3B total
parameters \cite{mobilemoe}. This is a property of the retrieved catalogue
record, not a reason to silently rewrite the original cohort. The release
retains every original match and puts the substantive distinctions in
separate, source-supported annotations.
